\documentclass[10pt]{article}

% Shared presentation layer for the standalone R0--R7 development documents.
% Method equations remain authoritative in the canonical idea sketch.
\usepackage[a4paper,margin=18mm,headheight=14pt]{geometry}
\usepackage{fontspec}
\usepackage{xeCJK}
\xeCJKsetup{CJKspace=true}
\setmainfont{DejaVu Serif}
\setsansfont{DejaVu Sans}
\setmonofont{DejaVu Sans Mono}
\setCJKmainfont{Noto Serif CJK KR}
\setCJKsansfont{Noto Sans CJK KR}
\setCJKmonofont{Noto Sans Mono CJK KR}

\usepackage{amsmath,amssymb}
\usepackage{booktabs,longtable,tabularx,array}
\usepackage{enumitem}
\usepackage{xcolor}
\usepackage{hyperref}
\usepackage{fancyhdr}
\usepackage{microtype}
\usepackage[most]{tcolorbox}

\definecolor{CoreBlue}{HTML}{1F4E79}
\definecolor{CoreLight}{HTML}{EAF3F8}
\definecolor{StopRed}{HTML}{9C0006}
\definecolor{StopLight}{HTML}{FCE8E6}
\definecolor{GateGreen}{HTML}{38761D}
\definecolor{GateLight}{HTML}{EAF4E2}
\definecolor{DecisionOrange}{HTML}{B45F06}
\definecolor{DecisionLight}{HTML}{FFF2CC}

\hypersetup{
  colorlinks=true,
  linkcolor=CoreBlue,
  urlcolor=CoreBlue
}

\setlength{\parindent}{0pt}
\setlength{\parskip}{0.3em}
\setlength{\emergencystretch}{3em}
\setlist[itemize]{leftmargin=1.45em,itemsep=0.1em,topsep=0.2em}
\setlist[enumerate]{leftmargin=1.75em,itemsep=0.14em,topsep=0.22em}
\renewcommand{\arraystretch}{1.15}
\newcolumntype{L}[1]{>{\raggedright\arraybackslash}p{#1}}
\newcommand{\code}[1]{{\small\nolinkurl{#1}}}
\newcommand{\must}[1]{\textcolor{CoreBlue}{\textbf{#1}}}
\newcommand{\haltitem}[1]{\textcolor{StopRed}{\textbf{#1}}}
\newcommand{\passitem}[1]{\textcolor{GateGreen}{\textbf{#1}}}
\newcommand{\decisionitem}[1]{\textcolor{DecisionOrange}{\textbf{#1}}}

\newenvironment{contractbox}[1]
  {\begin{tcolorbox}[colback=CoreLight,colframe=CoreBlue,title={#1},fonttitle=\bfseries,before upper=\raggedright]}
  {\end{tcolorbox}}
\newenvironment{decisionbox}[1]
  {\begin{tcolorbox}[colback=DecisionLight,colframe=DecisionOrange,title={#1},fonttitle=\bfseries,before upper=\raggedright]}
  {\end{tcolorbox}}
\newenvironment{stopbox}[1]
  {\begin{tcolorbox}[colback=StopLight,colframe=StopRed,title={#1},fonttitle=\bfseries,before upper=\raggedright]}
  {\end{tcolorbox}}


\hypersetup{
  pdftitle={R0 Contract and Schema --- Wind3DGS Development Specification},
  pdfauthor={Wind3DGS Project}
}

\pagestyle{fancy}
\fancyhf{}
\lhead{\small Wind3DGS Development --- R0}
\rhead{\small Contract and Schema}
\cfoot{\thepage}

\title{\textbf{R0 --- Contract and Schema}\\
개발 명세와 동결 체크리스트}
\author{Wind3DGS Project}
\date{2026-08-22; 상태 감사 2026-09-14}
\renewcommand{\DevChecklistDate}{2026-09-14}

\begin{document}
\maketitle

\begin{contractbox}{문서 권위와 사용법}
이 문서는 canonical idea sketch의 구현 파생 문서다. 방법 수식과 claim의 최종 권위는
\code{3dgs_response_distilled_global_local_wind_dynamics_2026-08-22.tex}에 있다.
특히 \code{eq:rd-setup-input}, \code{eq:rd-frame-input},
\code{eq:rd-typed-control-map}, \code{eq:rd-gaussian-token},
\code{eq:rd-response-package}, \code{eq:rd-area-mass},
\code{eq:rd-attachment-whitening}, \code{eq:rd-exact-zoh},
\code{eq:rd-wind-pullback}, \code{eq:rd-mean-transport}를 참조한다.
이 문서는 구현하면서 상태와 구체적 선택을 갱신하되, sketch의 계약을 조용히 바꾸지 않는다.
\end{contractbox}

\DevChecklistPage{R0}{Teacher 하위 계약 진행 · R0 미완료}{
\DevProgressRow{0}{\hyperref[sec:r0-chapter-1]{1. 목적과 소유권}}{대기}{범위와 owner를 executable contract 및 manifest에 연결}
\DevProgressRow{0}{\hyperref[sec:r0-chapter-2]{2. 진입 입력}}{대기}{Domain/metric/material/attachment seed와 revision 동결}
\DevProgressRow{0}{\hyperref[sec:r0-chapter-3]{3. 이미 동결된 계약}}{진행}{Teacher 하위 계약과 target schema의 구분·불변량 검증}
\DevProgressRow{0}{\hyperref[sec:r0-chapter-4]{4. 구현 체크리스트}}{진행}{Setup/frame schema, validator와 serialization 완성}
\DevProgressRow{0}{\hyperref[sec:r0-chapter-5]{5. Open Design Decisions}}{대기}{미결정 option과 freeze 시점·근거 확정}
\DevProgressRow{0}{\hyperref[sec:r0-chapter-6]{6. 산출물}}{대기}{Canonical registry와 stage artifact manifest 생성}
\DevProgressRow{0}{\hyperref[sec:r0-chapter-7]{7. 필수 fixture와 metric}}{대기}{Roundtrip/permutation/privilege/단위 fixture 통과}
\DevProgressRow{0}{\hyperref[sec:r0-chapter-8]{8. 종료 조건과 실패 경로}}{대기}{R0 종료 항목과 hand-off report 승인}
}{
\begin{contractbox}{현재 확인된 범위}Teacher registry 등 일부 하위 계약 구현은 있다. Target response schema, 선택 동결 및 canonical R0 fixture/report의 완료 근거는 아직 없다. R1의 개발 진단을 R0 전체 완료로 승계하지 않는다.\end{contractbox}
}

\begin{contractbox}{2026-09-14 하위 구현과 R0 완료의 구분}
Teacher 실행에서는 SI 물성, FP64 hi/lo 상태, 코드/입력 hash 동결, 단계별 report와
원시 checkpoint, 중단 prefix 보존 및 재생 검증을 구현했다. 이는 개발 run의 재현 계약이다.
Target setup/frame/response package의 machine-readable schema, coordinate/permutation,
physical-valid/rank/privilege 검증과 canonical serialization fixture 전체를 완료했다는 근거는 아니다.
따라서 R0의 미완료 체크는 유지한다. R1의 GPU/굽힘 결과를 R0 target 계약 완료로 승계하지 않는다.
구현 계약은 \code{code/docs/gravity_wrinkle_test.md}, 검증 근거는
\code{experiments/R1_teacher_velocity_reset/timestep_search/cloth_coarse/}의 결과 문서를 따른다.
\end{contractbox}

\tableofcontents
\newpage

\section{목적과 소유권}
\label{sec:r0-chapter-1}

R0의 목적은 이후 R1--R7이 같은 입력, 단위, tensor 의미, artifact와 실패 분모를 사용하도록
\textbf{실행 경계}를 먼저 고정하는 것이다. 이 단계에서는 teacher data나 network 성능을 만들지 않는다.

\begin{longtable}{L{0.24\linewidth}L{0.30\linewidth}L{0.36\linewidth}}
\toprule
소유자 & R0가 소유하는 것 & R0가 소유하지 않는 것\\
\midrule
\endfirsthead
\toprule
소유자 & R0가 소유하는 것 & R0가 소유하지 않는 것\\
\midrule
\endhead
Contract/schema & setup/frame API, tensor semantic axis, dtype/unit, version, validation, serialization/hash 경계 & teacher constitutive model, model depth, 최종 loss weight\\
Physics boundary & SI-only, sole mass owner, traction/evaluator identity, target privilege boundary & teacher discretization 및 수렴 결과, learned response 정확도\\
Data boundary & source-object group, target-allowed field, sealed test identity, failure denominator policy & 실제 K0/K1 sample 생성과 split 통계\\
Evaluation boundary & \code{MetricRegistry}와 \code{ThresholdRegistry}의 형식 및 freeze lifecycle & dev 결과를 보기 전 정할 수 없는 수치 threshold\\
Stage interface & R0--R7 artifact producer/consumer와 완료 보고 형식 & 각 후속 stage의 구현 내부 구조\\
\bottomrule
\end{longtable}

\must{R0 성공의 의미:} 같은 package를 serialize--deserialize해도 evaluator 결과가 동일하고,
Gaussian 입력 순서가 달라도 stable ID로 되돌린 결과가 같으며, teacher-only field가 target forward에
들어갈 수 없고, 모든 물리 scalar와 failure가 registry에 기계적으로 기록되는 상태다.

\section{진입 입력}
\label{sec:r0-chapter-2}

\subsection{권위 문서와 초기 상태}

\begin{itemize}
  \item Canonical sketch와 현행 R0--R7 checklist의 동일 revision/hash.
  \item Narrow core domain: flat-rest single-layer strip/rectangular flag, hard edge 또는 작은 declared patch
  attachment, known metric, 2--3 isotropic material preset, prescribed one-way wind, small-to-moderate deformation.
  \item 기존 explicit-scaffold 및 legacy code/experiment는 참고 구현 또는 baseline일 뿐 새 R0 완료 근거가 아니다.
  \item R0 진입 시 sealed test의 내용은 열지 않는다. Source-object group ID와 sealing hash만 만든다.
\end{itemize}

\subsection{Setup과 frame 의미 입력}

Setup 의미 입력은 sketch \code{eq:rd-setup-input}의
\[
  \mathcal I_{\mathrm{setup}}=
  (\mathcal G^0,\vartheta_{\mathrm{metric}},\vartheta_{\mathrm{mat}},
   \mathcal C_{\mathrm{attach}},\vartheta_{\mathrm{aero}},\vartheta_{\mathrm{domain}})
\]
이다. \(\mathcal G^0\)는 rest static GS다. Metric payload는
\[
  \vartheta_{\mathrm{metric}}=(L_0,A_{\mathrm{ref}},M_{\mathrm{ref}})>0
\]
이며 세 값 모두 positive다. \(\vartheta_{\mathrm{aero}}\)는 domain-wide frozen
\(\kappa_{\mathrm{aero}}\)[kg/m\(^3\)]와 \code{traction_law_id}만 물리 공력 항으로 소유한다.
\(\vartheta_{\mathrm{domain}}\)은 \(\tau_{\mathrm{guard}}\)[Pa]와 activation failure/OOD policy를 소유하며,
guard를 artist control이나 재료/공력계수로 노출하지 않는다.
% [논문 작성 참고 -- 수정 이유]
% kappa_aero=0.5*rho_air*C_n만 현재 힘에서 식별 가능하므로 rho_air와 C_n을 schema field로 분리하지 않는다.
% tau_guard는 결과를 맞추는 물리 자유도가 아니라 overflow 검출용 수치 계약이므로 DomainContract가 소유한다.
Frame 의미 입력은 sketch \code{eq:rd-frame-input}의
\[
  \mathcal I_t=(W_t,\vartheta_{\mathrm{ctrl},t},\Delta t,s_{t-1},\mathcal P),
  \qquad \vartheta_{\mathrm{ctrl},t}=(g_W(t),g_K,g_D)
\]
이다. Setup compiler는 package를 한 번 만들고, frame evaluator는 cached package와 명시적 state만 소비한다.

Target 입력에 mesh vertex/connectivity, geodesic, mesh--GS correspondence, teacher trajectory,
dynamic target RGB/video 또는 별도 opaque hidden state가 들어오면 R0 실패다.

\subsection{필수 선행 registry seed}

R0 시작 시 아직 값이 비어 있어도 다음 registry의 schema와 owner를 만든다.
\begin{itemize}
  \item \code{DomainContract}, \code{TypedControlContract}, \code{TeacherStudentVisibility}
  \item \code{TensorSchemaRegistry}, \code{CoordinateContract}, \code{SerializationContract}
  \item \code{MetricRegistry}, \code{ThresholdRegistry}, \code{StageArtifactRegistry}
  \item \code{ObjectGroupedSplit}, \code{FailureReasonRegistry}, \code{MethodIdentityRegistry}
  \item \code{RendererBranchContract}: base normal/covariance branch와 optional appearance branch 허용 여부;
  SH degree와 residual capacity는 R7이 optional branch를 열 때 dev-freeze
\end{itemize}

\section{이미 동결된 계약}
\label{sec:r0-chapter-3}

이 절은 설계 선택지가 아니다. 변경하려면 canonical sketch와 모든 downstream 문서를 함께 revision한다.

\subsection{물리 단위와 질량 소유권}

\begin{itemize}
  \item 외부/runtime contract는 SI-only다. \(L_0\)[m], \(A_{\mathrm{ref}}\)[m\(^2\)],
  \(M_{\mathrm{ref}}\)[kg], time/\(\Delta t\)[s], force[N], traction[Pa]를 사용한다.
  \item \(\Omega\)는 positive angular frequency[rad/s], \(Z=\operatorname{diag}(\zeta_j)\)는
  positive dimensionless modal damping ratio다. Cycles/s나 숨은 normalized time을 허용하지 않는다.
  \item \(M_{\mathrm{ref}}\)가 total mass의 유일한 owner다. Homogeneous V0에서
  \(m_i=(M_{\mathrm{ref}}/A_{\mathrm{ref}})a_i\)이며 independent mass head는 없다
  (sketch \code{eq:rd-area-mass}). Material preset은 stiffness/damping만 소유한다.
  \item Encoder normalized coordinate는 numerical feature일 뿐 별도 물리 unit mode가 아니다.
  Loss, validator와 evaluator 전에 SI payload로 복원한다.
\end{itemize}

\subsection{출력 의미와 runtime 경계}

Teacher-only schema는 element/DOF와 positive quadrature, mass-law 및 map-law/version을 구분한다.
고차 Teacher의 consistent mass와 signed shape map은 R1 및 sketch의
\code{eq:rd-teacher-consistent-mass}, \code{eq:rd-common-probe-map}을 따른다.
Student GS의 area-derived diagonal mass 또는 기존 convex map schema로 고차 payload를 해석하지 않는다.
최종 고차 registry/validator는 아직 미구현이며 native Teacher 하위 계약의 존재로 동결을 대체하지 않는다.

\begin{itemize}
  \item Setup network 출력은 absolute deformation이나 next pose가 아니라
  \textbf{causal force-to-motion operator의 parameter를 담은 response package}다
  (sketch \code{eq:rd-response-package}).
  \item Frame evaluator가 current wind와 package로 generalized response/state increment를 계산한다.
  Setup package와 frame increment를 같은 ``response output''으로 serialize하지 않는다.
  \item Global은 always-on passive package다. Local은 frozen Global이 설명한 component를 뺀 residual이며,
  Local budget/상태는 R5--R6 문서가 소유한다.
  \item Runtime package와 state는 rest GS 및 모든 target-allowed physical metadata의 identity/hash를 소유한다.
  Frame evaluator가 teacher mesh나 외부 hidden cache를 조회하지 않는다.
\end{itemize}

\subsection{Typed control과 evaluator identity}

R1의 checkpoint/velocity-reset은 teacher-only 탐색 기록이다. 이를 target frame input이나 기존 typed
artist control에 추가한 것으로 해석하지 않는다. 비정지/window 시작을 student에 사용할 때는
\(q,\dot q\) 초기화와 초기 재구성 오차의 명시적 계약이 필요하며 구체적인 알고리즘은 아직 미동결이다.

\begin{itemize}
  \item \(0\le g_W(t)\le g_W^+\), \(W_t^{\mathrm{eff}}=g_W(t)W_t\)이며 wind direction은 \(W_t\)가 소유한다.
  \item Package의 \(\Omega_S^{\mathrm{base}},Z_S^{\mathrm{base}}\)는 immutable이다. Evaluator만
  \(\Omega_S^{\mathrm{eff}}=\sqrt{g_K}\Omega_S^{\mathrm{base}}\),
  \(Z_S^{\mathrm{eff}}=g_DZ_S^{\mathrm{base}}\)를 파생한다
  (sketch \code{eq:rd-typed-control-map}).
  \item \(g_K,g_D>0\)는 rollout 동안 고정한다. 이를 바꾸면 derived-evaluator hash를 바꾸고 rest state에서
  재시작한다. Material/attachment 변경은 setup rerun과 새 package를 요구한다.
  \item Traction law는 \code{traction_law_id}=\code{two_sided_normal_quadratic_v1}이며,
  teacher와 target GS는 동일한 domain-wide \(\kappa_{\mathrm{aero}}\)를 사용한다.
  % [논문 작성 참고 -- 수정 이유]
  % 공기밀도/법선계수/풍속 gain을 동시에 자유롭게 두면 F가 kappa_aero*g_W^2에 비례해 식별이 닫히지 않는다.
  % 그래서 kappa_aero는 package identity에 고정하고 runtime에서는 g_W만 바람 세기 control로 허용한다.
  \item Numerical guard는 \code{traction_guard_id}=\code{vector_norm_before_area_v1}다.
  Vector-norm traction guard를 area 곱보다 먼저 적용하고 정량 core의 activation count는 0이어야 한다
  (sketch \code{eq:rd-wind-pullback}). Activation이 생긴 frame은 OOD/failure denominator에 남긴다.
  % [논문 작성 참고 -- 수정 이유]
  % guard activation을 허용하면 quadratic law가 clipped law로 바뀌므로 정상 결과의 안정화 장치로 쓰지 않는다.
  \item Evaluator identity는 \code{integrator_id}=\code{exact_zoh_augexp_v1},
  \code{force_sample_time}=\code{frame_start_v1}다.
  Frame-start load를 한 frame hold하고 mode별 exact ZOH를 한 번만 적용한다
  (sketch \code{eq:rd-exact-zoh}).
  \item Zero ambient \(W_t=0\)은 relative-velocity drag를 남길 수 있다. Strict passive free decay는
  aerodynamic evaluation을 끈 \(Q=0\) fixture로 별도 표기한다.
\end{itemize}

\subsection{Topology, split와 permutation 원칙}

\begin{itemize}
  \item Target inference는 persistent physical adjacency graph를 생성, 저장 또는 소비하지 않는다.
  Patch/attention relation은 같은 network forward 안에서만 사용하는 transient computation이다.
  \item Mesh topology와 correspondence는 teacher simulation, training loss와 evaluation remap에만 허용된다.
  \item 같은 source object의 모든 mesh, independent GS variant, material, wind와 frame은 같은 split에 둔다.
  원본 trajectory, checkpoint 분기, intervention 시각과 파생 window/patch의 lineage를 보존하고
  velocity-reset 분기도 원본과 같은 source-object group에 둔다.
  \item Gaussian-to-token aggregation은 set 입력을 읽고 token-to-Gaussian query decoder는 입력 permutation에
  equivariant해야 한다. Serialized 순서가 model 의미를 바꾸어서는 안 된다.
  \item Test normalization, threshold, rank, budget와 loss weight는 test open 뒤 바꾸지 않는다.
\end{itemize}

\subsection{Package 불변량}

\begin{itemize}
  \item Area는 nonnegative이고 declared normalization domain에서 합이 \(A_{\mathrm{ref}}\)다.
  Mass는 area에서 파생되고 합이 \(M_{\mathrm{ref}}\)다. Physical-valid set의 정확한 normalization domain은
  아직 Open Design Decision이다.
  \item Attachment row는 exact zero projection 후 mass whitening하며 \(B^\top M_xB=I\)를 만족한다
  (sketch \code{eq:rd-attachment-whitening}). Rank floor 아래 Gram은 축소 또는 package reject 대상이다.
  어느 정책을 택할지는 아직 동결하지 않았다.
  \item Normal/covariance branch는 angular field \(B_R\) 또는 versioned local-affine update 중 하나만
  전체 domain에 사용한다. 어느 branch인지는 R0 종료 전에 동결해야 한다.
  \item Base mean/covariance transport는 attachment zero, rest identity, proper rotation, SPD/finite covariance를
  만족해야 한다 (sketch \code{eq:rd-mean-transport}).
\end{itemize}

\section{구현 체크리스트}
\label{sec:r0-chapter-4}

\subsection{Schema와 API}

\begin{itemize}
  \item[$\square$] \code{response_contract_v1}에 setup compiler, package validator, frame evaluator,
  state reset과 serialization API를 명시한다.
  \item[$\square$] 함수별 allowed/forbidden argument, return type, device/dtype policy와 error type을 등록한다.
  \item[$\square$] Setup input, response package, runtime state와 frame trace를 서로 다른 type/version으로 둔다.
  \item[$\square$] Variable \(N\) 입력 batching, padding/mask, ragged Local support와 batch axis 규칙을 기록한다.
  \item[$\square$] 모든 tensor에 semantic axis, physical unit, dtype, finite/positivity condition, stable-ID relation을 붙인다.
  \item[$\square$] Package rank/condition, law/evaluator ID, model/config/control-range와 rest-GS hash를 validator가 확인한다.
  \item[$\square$] Material/attachment/control 변경 시 package 또는 derived-evaluator hash/reset 규칙을 자동 검사한다.
\end{itemize}

현재 schema가 반드시 표현해야 할 semantic tensor는 다음과 같다. 아래 axis는 의미 요구사항이며,
contiguous layout, packed covariance 및 ragged encoding은 Open Design Decision이다.

\begin{longtable}{L{0.24\linewidth}L{0.28\linewidth}L{0.38\linewidth}}
\toprule
Payload & 필수 semantic axis & 필수 validator\\
\midrule
\endfirsthead
\toprule
Payload & 필수 semantic axis & 필수 validator\\
\midrule
\endhead
Rest GS & Gaussian \(N\), xyz, covariance, opacity/allowed appearance & stable ID unique, covariance SPD/finite, rest hash 일치\\
Metric/material & scalar 또는 preset axis & SI unit, positivity/range, sole-mass-owner 위반 없음\\
Attachment/validity & Gaussian \(N\) 또는 geometric-region description & declared owner/version, deterministic materialization\\
Area/mass/normal & \(a,m,\gamma_{\rm valid},\gamma_{\rm ori}\) on \(N\); \(n^0\) on \(N\times3\) & nonnegative/finite, unit normal on accepted rows, sum/zero policy\\
Global field & \(B_G\) on \(N\times3\times r_G\); \(\Omega_G,Z_G\) on \(r_G\) & attached/invalid row, rank, Gram identity, positive pole/damping\\
Optional angular field & \(B_{R,G}\) on \(N\times3\times r_G\) 및 Local 대응 axis & branch ID, shared coordinate, unit/finite condition\\
Local package & slot \(K_L\), ordered support, support xyz, Local rank & support uniqueness/order, attachment/validity, sparse bounds\\
Runtime state & batch, Global/Local coordinate와 velocity, Local state/timer/fade & package identity, finite state, reset provenance\\
Frame trace & frame/time, control, load, state transition, timing, failure & monotone time, actual \(\Delta t\), method identity와 denominator key\\
\bottomrule
\end{longtable}

\subsection{Coordinate와 permutation contract}

\begin{itemize}
  \item[$\square$] World, object-canonical, normalized-feature와 generalized coordinate를 이름과 transform으로 분리한다.
  \item[$\square$] Object center \(o\), positive \(L_0\), attachment cue, orientation/sign confidence와 inverse transform을 저장한다.
  \item[$\square$] Position은 \((\mu^0-o)/L_0\), covariance는 \(\Sigma^0/L_0^2\)로만 feature normalize하고
  physical output은 SI로 복원한다 (sketch \code{eq:rd-gaussian-token}).
  \item[$\square$] Random input permutation 전후 package를 stable ID로 gather했을 때 area/normal/field와 rollout이 일치한다.
  \item[$\square$] Serialization order, model compute order와 renderer order를 분리하고 모든 gather/scatter에 stable ID를 요구한다.
  \item[$\square$] Proper rigid transform, uniform metric scale, symmetry/sign ambiguity와 rejection behavior를 fixture로 만든다.
\end{itemize}

\subsection{Physical-valid set와 package validation}

\begin{itemize}
  \item[$\square$] \code{PhysicalValidPolicy}에 valid-set owner와 hard/soft field를 기록한다.
  Area normalization domain, invalid \(a,m,F,B,n\) 처리와 asset reject rule도 같은 contract에 명시한다.
  \item[$\square$] Floater/background/close two-sided layer가 load나 mode mass를 몰래 소유하지 않는지 검사한다.
  \item[$\square$] Hard attachment와 physical-invalid를 별도 reason code로 유지하고 denominator에서 숨기지 않는다.
  \item[$\square$] Valid-set 변경이 area/mass/package/map hash를 모두 무효화하도록 dependency를 등록한다.
  \item[$\square$] Fixed-rank 미달, invalid covariance, orientation confidence 부족, nonfinite tensor와 OOD를
  predeclared package failure로 serialize한다.
\end{itemize}

\subsection{MetricRegistry와 ThresholdRegistry}

\begin{longtable}{L{0.24\linewidth}L{0.66\linewidth}}
\toprule
Registry & 필수 field\\
\midrule
\endfirsthead
\toprule
Registry & 필수 field\\
\midrule
\endhead
\code{MetricRegistry} & \code{metric_id/version}, 물리 quantity와 unit, probe/map/common-valid-mask ID,
normalization/reference scale, sample axis, time/sequence reduction, source-object outer aggregation,
seed aggregation, better direction, NaN/failure encoding, implementation/config hash\\
\code{ThresholdRegistry} & \code{threshold_id/version}, 연결 metric, 비교 연산자, 숫자와 unit,
absolute/relative 여부, 적용 domain/stage/Gate, tolerance source, dev freeze 시점, sealed-test 사용 금지,
override 금지와 registry hash\\
\code{FailureReasonRegistry} & OOD pre-exclusion과 in-envelope method failure를 분리하는 stable reason,
denominator inclusion, owning stage, first-failure precedence와 trace/report mapping\\
\end{longtable}

\begin{itemize}
  \item[$\square$] Metric 함수가 값뿐 아니라 denominator count, excluded/failure count와 registry hash를 반환한다.
  \item[$\square$] Gate는 자연어 ``의미 있는 개선''이 아니라 threshold ID들의 AND/OR 식을 manifest에 저장한다.
  \item[$\square$] Source object를 outer cluster로 두고 그 안의 GS/wind를 nested paired sample로 집계할 수 있게 한다.
  \item[$\square$] Test open 후 새 metric/threshold version을 같은 result table에 조용히 덮어쓰지 못하게 한다.
\end{itemize}

\subsection{Serialization과 reproducibility}

\begin{itemize}
  \item[$\square$] Schema version, endianness/dtype, tensor shape, unit, optional field presence와 dependency hash를 저장한다.
  \item[$\square$] Serialize--deserialize--serialize canonical bytes 또는 semantic hash의 round-trip identity를 검사한다.
  \item[$\square$] Unknown required field/version은 hard fail하고 unknown optional field 처리 규칙을 명시한다.
  \item[$\square$] Package/run에 source commit, config, checkpoint, hardware/precision, command와 replay seed를 저장한다.
  \item[$\square$] Secret, 개인 절대 경로와 teacher-only raw payload가 target package에 들어가지 않는지 검사한다.
\end{itemize}

\subsection{Stage artifact matrix}

다음 표는 각 문서가 소유할 최소 산출물과 소비 경계를 정의한다. 실제 파일명/포맷은
\code{StageArtifactRegistry}에서 동결한다.

\begin{longtable}{L{0.09\linewidth}L{0.30\linewidth}L{0.24\linewidth}L{0.26\linewidth}}
\toprule
Stage & Produced artifact & Primary consumer & 금지된 대체 근거\\
\midrule
\endfirsthead
\toprule
Stage & Produced artifact & Primary consumer & 금지된 대체 근거\\
\midrule
\endhead
R0 & named contracts/registries, \code{ObjectGroupedSplit}, \code{RendererBranchContract}, \code{R0ContractReport} & R1--R7 전부 & legacy milestone 또는 자연어-only 설정\\
R1 & \code{TeacherPackage}, \code{CanonicalProbeMap}, \code{StaticGSPackage}, \code{OracleGlobalPackage}, \code{BaseTransportReport}, \code{TeacherSpectrumStatus} & R2--R7 & network 결과로 reference 오류를 덮는 것\\
T1 & \code{T1InitCheckpoint}, \code{T1InitReport} & R2a & total measure 합만으로 initialization 합격\\
R2a & \code{R2aOverfitCheckpoint}, single-case fixed-seconds package/trace와 overfit report & R2b & short one-step loss만으로 합격\\
R2b & \code{R2GlobalCheckpoint}, \code{R2bHandoff}, \code{GateATPrerequisiteReport} & R3--R5 & best schedule의 사후 선택\\
R3 & \code{MultiResplatGroupManifest}, measure/response reports, \code{GateAReport}, \code{GlobalConfigSet}, \code{R3RouteConfig}, \code{R3RouteCheckpoint} & R4, claim registry & construction상 area/mass 합만 증거로 사용\\
R4 & \code{DataSplitManifest}, \code{BestGlobalCheckpoint}, \code{GlobalResponsePackageSet}, \code{GateBReport}, \code{R5GlobalHandoff} & R5, R7 & model checkpoint와 per-asset package를 한 artifact로 합침\\
R5 & \code{GlobalResidualAtlas}, \code{AlwaysOnLocalCheckpoint}, \code{AlwaysOnLocalPackageSet}, \code{R5ProfilerFixture}, \code{GateCReport} & R6, R7 & teacher/test residual로 support 사후 선택\\
R6 & \code{ConditionalRuntimeConfig}, \code{ConditionalRuntimeTrace}, \code{GateDReport} & R7 & idealized FLOPs만으로 speed claim\\
R7 & \code{RendererEvidenceContract}, \code{ChampionFailureRegistry}, \code{PaperEvidencePackage} & paper & appearance로 실패한 motion Gate 구제\\
\bottomrule
\end{longtable}

위 이름은 \code{StageArtifactRegistry}의 canonical artifact ID다. 실제 파일명은 달라도 producer와 consumer가
같은 ID와 content hash를 사용해야 하며 자연어 alias로 대체하지 않는다. 특히 학습된 model checkpoint와
그 model이 declared asset마다 compile한 \code{ResponsePackage} set은 서로 다른 artifact로 직렬화하고
dependency hash로만 연결한다.

\section{Open Design Decisions}
\label{sec:r0-chapter-5}

\begin{decisionbox}{해석 규칙}
아래 ``권장 기본값''은 구현을 시작하기 위한 제안이며 아직 동결된 방법이 아니다.
실제 선택은 지정된 시점에 decision record, config/schema version과 fixture 결과를 함께 남긴 뒤에만 계약이 된다.
\end{decisionbox}

\begin{longtable}{L{0.19\linewidth}L{0.25\linewidth}L{0.30\linewidth}L{0.16\linewidth}}
\toprule
결정 항목 & 허용 선택지 & 권장 기본값(미동결) & 동결 시점\\
\midrule
\endfirsthead
\toprule
결정 항목 & 허용 선택지 & 권장 기본값(미동결) & 동결 시점\\
\midrule
\endhead
Physical-valid 정책 & Gaussian hard valid set / asset-level reject / 둘의 사전 선언 조합 & hard valid set에서만 area를 normalize하고 invalid \(a,m,F,B\)를 exact zero; confidence가 전체적으로 낮으면 asset reject & R0 package schema 종료 전\\
Canonical frame & attachment 축+geometry cue / sign-invariant feature+confidence reject & attachment-defined primary cue, symmetry에서는 sign-invariant feature와 reject reason & Accepted R1 data 생성 전\\
Tensor layout & dense padded / ragged offsets / 혼합 & Global dense, Local은 ordered index+offset ragged, batch mask 명시 & 첫 package serializer 전\\
Stable ordering & input order 유지+ID gather / canonical stable-ID sort & serialization은 stable-ID sort, network는 임의 permutation을 허용하고 output을 ID로 scatter & R0 permutation fixture 전\\
Rank 부족 & rank 축소 / fixed-rank package reject & V0는 fixed rank 미달 package reject로 contract를 단순화 & R2a dataset freeze 전\\
Normal/\allowbreak covariance & shared-coordinate angular field \(B_R\) / versioned local-affine mean update & persistent graph가 필요 없는 \(B_R\); supervision과 package 크기가 과도하면 local-affine & R1 traction/transport fixture 전\\
Angular composition & \(B_R\) branch에서 Global/Local 합산 후 한 번의 exp-map / 사전 선언한 다른 proper composition & deterministic order와 proper \(\mathrm{SO}(3)\), attachment identity를 동시에 만족하는 가장 단순한 규칙 & R1 traction/transport fixture 전\\
Covariance encoding & full \(3\times3\) / symmetric packed-6 / eigensystem & packed-6 저장+validator에서 full SPD 복원 & serializer v1 동결 전\\
Serialization & safe tensor container+JSON manifest / schema binary format & 비실행형 tensor container와 canonical JSON manifest를 분리하고 semantic content hash 사용 & Canonical R1 artifact 생성 전\\
Metric normalization & SI absolute / reference-scale dimensionless / 둘 다 보고 & Gate primary는 predeclared dimensionless, SI absolute는 diagnostic으로 병기 & R2 loss/metric freeze 전\\
Threshold lifecycle & 모든 수치 R0 고정 / stage별 dev freeze & schema는 R0, 숫자는 해당 stage dev에서 test open 전에 freeze하고 immutable version 발행 & 각 stage test open 전\\
\bottomrule
\end{longtable}

Physical-valid 정책과 normal/covariance branch는 area, load, current normal과 method identity를 바꾸는
결정이다. 두 항목과 angular composition을 동결할 때는 R0 decision record에만 두지 않고 canonical
sketch와 master roadmap의 semantics도 먼저 갱신한다. 그 revision/hash가 일치해야 R0를 종료할 수 있다.

\section{산출물}
\label{sec:r0-chapter-6}

R0 완료 시 최소한 다음 canonical artifact ID가 존재해야 한다. 괄호 안 파일명은 예시이며 registry가 실제 경로를 소유한다.

\begin{itemize}
  \item \code{domain_contract.json}, \code{typed_control_contract.json}\\
  \code{teacher_student_visibility.json}, \code{object_grouped_split.json}
  \item \code{tensor_schema_registry.json}, \code{coordinate_contract.json}\\
  \code{serialization_contract.json}, \code{physical_valid_policy.json}
  \item \code{metric_registry.json}, \code{threshold_registry.json}\\
  \code{failure_reason_registry.json}, \code{method_identity_registry.json}
  \item \code{stage_artifact_registry.json}, \code{response_contract_v1}의 machine-readable schema와 API 문서
  \item \code{RendererBranchContract}: base normal/covariance branch와 optional appearance 허용 flag
  \item Minimal valid/invalid package fixture, round-trip package, frame input/state/trace fixture
  \item \code{R0ContractReport}: 결정 기록, 미결정 항목, fixture command/result,
  source/config/schema hash와 다음 stage 허용 여부
\end{itemize}

Package에는 immutable rest GS identity, target metric/material/attachment/aero/domain identity,
physical-valid/normal branch ID, response field와 pole 의미, rank/condition, law/evaluator ID와 model/config hash가
닫힌 형태로 들어가야 한다. Runtime state와 frame trace는 package 파일과 분리한다.
% [논문 작성 참고 -- 수정 이유]
% aero identity에는 kappa_aero와 normal-only law만, domain identity에는 tau_guard와 activation policy만 넣어
% 논문에서 물리 파라미터와 numerical safeguard를 분리해 보고할 수 있게 한다.

\section{필수 fixture와 metric}
\label{sec:r0-chapter-7}

\begin{longtable}{L{0.22\linewidth}L{0.50\linewidth}L{0.18\linewidth}}
\toprule
Fixture & 합격 조건 & 실패 owner\\
\midrule
\endfirsthead
\toprule
Fixture & 합격 조건 & 실패 owner\\
\midrule
\endhead
Schema round trip & 모든 required tensor/unit/ID/hash가 semantic identity로 복원되고 evaluator 결과가 bitwise 또는 declared tolerance 안에서 동일 & serialization 또는 schema\\
Unit dimensionality & SI sample의 area--mass--force--generalized load--state unit이 validator와 analytic fixture에 일치 & domain/schema\\
Permutation & 여러 random Gaussian permutation 후 stable-ID gather 결과와 rollout이 동일 & model API/schema\\
Coordinate & rest identity, proper rigid transform, uniform scale와 inverse transform이 declared tolerance 통과 & coordinate contract\\
Hash invalidation & rest GS, material, attachment, valid set, law, control range 또는 schema를 바꾸면 올바른 downstream hash가 달라짐 & artifact registry\\
No leakage & target forward signature, serialized package와 batch에 teacher-only field가 없음; object-group split 교차 없음 & visibility/split\\
Valid set & invalid Gaussian의 area/mass/load/field ownership과 denominator가 frozen policy와 정확히 일치 & valid policy\\
Package invariant & positivity, attachment zero, Gram identity, rank/condition, branch ID와 rest hash 통과 & package validator\\
Control/reset & \(g_K,g_D\), material, attachment 변경이 각각 정해진 derived hash/package rerun/reset을 유발 & API/state\\
Failure denominator & NaN, invalid covariance, rank fail, rollout fail과 OOD가 서로 다른 stable reason으로 집계됨 & metric/failure registry\\
\bottomrule
\end{longtable}

모든 fixture는 deterministic seed, command, source/config hash와 측정 tolerance를 남긴다.
수동 screenshot이나 log 문구만으로 합격 처리하지 않는다.

\section{종료 조건과 실패 경로}
\label{sec:r0-chapter-8}

\subsection{R0 종료 조건}

다음이 모두 참일 때만 canonical R1 acceptance와 학습 label 승격을 진행한다.
격리된 teacher-only 개발 진단·샘플 생성은 허용하지만 \code{training_eligible=false}로 유지하며,
그 하위 검사 결과를 R0 전체 완료로 승계하지 않는다.
\begin{itemize}
  \item[$\square$] setup/frame/package/state/trace type과 tensor semantic axis가 machine-readable schema로 닫혔다.
  \item[$\square$] Coordinate, physical-valid, rank, normal/covariance와 serialization 선택이 decision record로 동결됐다.
  \item[$\square$] SI/mass owner, target privilege, permutation, package invariant와 reset fixture가 모두 통과했다.
  \item[$\square$] Metric/threshold/failure registry의 형식, stage artifact matrix와 sealed-test lifecycle이 동결됐다.
  \item[$\square$] R0 completion report가 미결정 수치의 owner와 정확한 동결 시점을 명시한다.
\end{itemize}

\begin{stopbox}{R0 실패와 routing}
Setup package와 frame response increment를 혼용하거나, topology/correspondence leak를 닫지 못하거나,
unit/mass/valid-set owner가 두 개 이상이거나, permutation/serialization 후 의미가 달라지면
accepted Teacher 생성과 학습 경로 진입을 진행하지 않는다. Schema와 fixture owner를 고쳐 같은 R0를 다시 실행한다.
Dataset 확대, Local module 또는 renderer는 이 실패를 해결하는 경로가 아니다.
\end{stopbox}

R0에서 아직 숫자를 정할 근거가 없는 threshold는 임의값으로 ``완료'' 처리하지 않는다.
대신 registry field, dev 산출물 owner, 동결 deadline과 test-access 금지를 닫으면 된다.

\end{document}
